\section{Experimental Design}

\subsection{Dataset}

We evaluate on the Cataract-Seg dataset~\cite{cataractseg2023roboflow}, a publicly available collection of anterior segment images annotated for pixel-wise segmentation of cataract-related structures including the cornea, pupil, and lens. The dataset provides pixel-level masks suitable for binary and multi-class segmentation evaluation. This domain is closely related to the surgical scene annotations established in CaDIS~\cite{grammatikopoulou2021cadis}, but focuses on still clinical images rather than intraoperative video frames.

The dataset is partitioned into training, validation, and test sets following a 70\%/15\%/15\% stratified split, yielding 210 training, 45 validation, and 45 test images. No patient-level overlap exists across splits.

\subsection{Preprocessing}

All images are resized to $512 \times 512$ pixels and RGB values normalized to $[0, 1]$. Edge maps are computed per image using Canny filtering applied to the grayscale conversion, with thresholds $t_1 = 50$ and $t_2 = 150$, and replicated to three channels before input to Encoder B. The grayscale conversion uses the standard luminosity formula $Y = 0.299R + 0.587G + 0.114B$, which weights green most heavily as it contributes most to perceived brightness. This preprocessing is applied at both training and test time. Preprocessing is applied identically across all baseline and proposed models to ensure fair comparison.

\subsection{Data Augmentation}

During training, the following augmentations are applied on-the-fly to both the RGB image and its corresponding mask: random horizontal flip (p = 0.5), random rotation within $\pm 15^{\circ}$, and random brightness and contrast jitter (factor range $[0.8, 1.2]$). Edge maps are recomputed from the augmented RGB image after geometric transforms. No augmentation is applied at validation or test time.

\subsection{Baselines}

Four models are evaluated to isolate the contribution of cross-attention fusion:

\begin{table}[h]
\centering
\caption{Models compared in the evaluation.}
\label{tab:baselines}
\begin{tabular}{ll}
\hline
\textbf{Model} & \textbf{Description} \\
\hline
U-Net (RGB) & Standard U-Net, RGB input only \\
U-Net (Edge) & Standard U-Net, edge map input only \\
U-Net (Early Fusion) & Single encoder, RGB+Edge concatenated at input \\
Proposed & Dual encoder + bottleneck cross-attention \\
\hline
\end{tabular}
\end{table}

All models use the same U-Net backbone, optimizer, and training schedule to ensure fair comparison. The first three baselines share the proposed model's decoder and loss function.

\subsection{Metrics}

Model performance is reported using Intersection over Union (IoU), Dice coefficient, and F1 score. For binary segmentation, let $\hat{M}$ be the predicted mask and $M$ the ground-truth mask. IoU measures the overlap between prediction and ground truth divided by their total area:
%
\begin{equation}
\text{IoU} = \frac{|\hat{M} \cap M|}{|\hat{M} \cup M|} = \frac{\text{overlap}}{\text{union}}
\end{equation}
%
A value of 1 means perfect overlap, 0 means no overlap. The Dice coefficient is a related measure that weights overlap more heavily:
%
\begin{equation}
\text{Dice} = \frac{2|\hat{M} \cap M|}{|\hat{M}| + |M|}
\end{equation}
%
Both metrics penalize false positives (predicting a pixel that is not in the ground truth) and false negatives (missing a pixel that is in the ground truth). IoU is the primary metric, as it is more conservative (IoU $\leq$ Dice) and standard in segmentation benchmarks. F1 is mathematically equivalent to Dice in binary segmentation and reported in the same column.

\subsection{Implementation Details}

Models are implemented in PyTorch using \texttt{segmentation-models-pytorch} for the U-Net backbone. Cross-attention is implemented with \texttt{torch.nn.MultiheadAttention} ($n_\text{heads} = 8$). All models are trained with AdamW (lr $= 10^{-4}$, weight decay $= 10^{-2}$) for 100 epochs with a batch size of 8. AdamW decouples weight decay from the gradient update, which improves generalization compared to standard Adam. The learning rate was chosen based on preliminary experiments and kept constant throughout training without a scheduler, as the loss curves showed smooth convergence without plateauing. The combined BCE+Dice loss is used throughout. BCE provides strong pixel-level gradients that help the model converge quickly, while Dice loss addresses the class imbalance between foreground and background regions by directly optimizing the overlap between prediction and ground truth. The two losses are summed without weighting, as this combination has been shown to be effective in medical segmentation tasks without requiring additional hyperparameter tuning. Metrics are computed with \texttt{torchmetrics}.

The proposed dual-encoder model has approximately 19.6 million parameters, compared to 15.0 million for a standard U-Net. This 30\% increase in parameters is modest relative to the architectural complexity added: two independent encoders, two cross-attention modules, and four skip projection layers. Inference time increases by approximately 40\% over a single U-Net, from 8 ms to 11 ms per image on an A100 GPU, which remains suitable for real-time clinical applications. The edge encoder, cross-attention modules, and skip projections together account for the additional parameters and computation. Experiments were run on a single NVIDIA A100 GPU (40 GB) via Google Colab. The random seed is fixed to 42 for all runs; code is available at \url{https://github.com/DaniloDuque/multimodal-cataract-seg}.
